\chapter{Threats to Validity}
\label{chap:validity}

Threats follow the construct, internal, external, and conclusion validity
categories of Wohlin et al.\ \cite{wohlin2012experimentation}. The intrinsic data
and design limitations are summarized in
Section~\ref{sec:method-limitations}.

\section{Construct Validity}
\label{sec:validity-construct}

Repository inactivity is the inverse of the future-window maintenance rule
\((t, t + 12\text{ months}]\)
(Section~\ref{sec:method-definitions}). It is a proxy: stable projects can need
little activity without being abandoned, and development may move to other
platforms or release channels. The label therefore cannot prove abandonment.
The two-of-four thresholds are heuristic, and merged pull requests proxy
maintainer interaction. Requiring sustained activity reduces incidental labels
but does not remove this threat.

The sensitivity analysis re-labels all \(50{,}000\) snapshots with one-of-four
and three-of-four variants (Section~\ref{sec:results-label-sensitivity}). The
inactive rate ranges from \(0.581\) to \(0.746\), while ROC-AUC remains within
\(0.958\)--\(0.961\). The result is not specific to the two-of-four threshold,
but the analysis cannot establish that public activity is the correct maintenance
construct. Individual cut-off sweeps remain future work
(Chapter~\ref{chap:conclusion}). The bot-filter ablation additionally shows that
excluding bots removes a one-directional bias that masks inactivity
(Appendix~\ref{appendix:bot-ablation}).

\section{Internal Validity}
\label{sec:validity-internal}

Temporal leakage is controlled by using timestamps \(\leq t\) for features, the
future window for labels, leaving incomplete windows unlabeled, and restricting
live enrichment to stable metadata
(Section~\ref{sec:method-leakage-prevention}). Activity persistence remains a
separate threat. A single recency signal nearly matches the model on the full
test set because many already-quiet repositories stay quiet
(Section~\ref{sec:results-baselines}). On active-at-\(t\), the model exceeds the
best trivial rule by \(0.061\) ROC-AUC
(Section~\ref{sec:results-slices}). The predictive contribution is claimed from
this subset, not the aggregate metric.

Quarterly observations create autocorrelated snapshots. Equal cohort sizes make
the test partition the final date, so every test repository also occurs earlier
in training (Section~\ref{sec:method-observation-window}). A
repository-disjoint split confines each repository to one partition and refits
the model on it (Appendix~\ref{sec:impl-grouped-split}). ROC-AUC decreases only
marginally, bounding optimism from recurrence and indicating discrimination on
repositories held out from the same frame
(Section~\ref{sec:results-robustness}).

The split is chronological but not embargoed
(Section~\ref{sec:method-label-availability}). Training and validation labels
close after the test observation date, and the boundary between training and
validation falls inside the \texttt{2024-10-01} cohort. Neither admits
information about a test row into its own features or label, so the reported
metric is not inflated by leakage across the boundary at \(t\). What they do
restrict is the interpretation: the design is a retrospective chronological
evaluation with overlapping outcome windows, so it does not demonstrate that the
model could have been fitted at the simulated prediction date.

\section{External Validity}
\label{sec:validity-external}

All labels and features come from GitHub events
(Section~\ref{sec:intro-scope}). The results do not cover GitLab, Codeberg,
self-hosted forges, or workflows based on mailing lists, such as the Linux
kernel. Platform choice may covary with project age, size, governance, and
maintenance, so the direction of this bias is unknown.\kiref{Claude Opus 5/1}
GitHub samples are also unrepresentative of software development generally: most
hosted repositories are personal and inactive, and roughly half of registered
users have no public activity
\cite{kalliamvakou2016perils}.\kiref{Claude Opus 5/10}

The foundation seed further requires a license and at least 100 stars. Results
do not generalize to very small, private, new, or low-visibility repositories;
the ``low popularity'' slice is relative to this frame
(Section~\ref{sec:results-slices}). Oversampling dormant and archived projects
also makes probabilities and precision conditional on sampled prevalence.
Deployment requires recalibration (Section~\ref{sec:method-limitations}).

The frame is also sampled retrospectively, which is the more serious of the two
sampling threats (Section~\ref{sec:method-retrospective-frame}). The seed was
drawn on \texttt{2026-06-11} using push recency and archive state at that date,
so bucket membership is decided by behavior falling inside the outcome windows
being predicted. The dormant bucket is inactive at a rate of \(1.000\) in the
test cohort, and roughly half of each cohort comes from buckets whose defining
condition largely fixes the outcome. Aggregate discrimination is consequently
measured on a population that is easier to separate than one sampled at \(t\),
and the aggregate figures should be read as conditional on this frame rather than
as an estimate of deployment performance. The active-at-\(t\) subset is less
exposed, since it excludes the already-quiet rows through which the buckets are
most legible, but it is drawn from the same frame and inherits the restriction.
Correcting this requires a frame built from repository state at each observation
date (Chapter~\ref{chap:conclusion}), which the reported run does not
do.\kiref{Claude Opus 5/11}

Finally, the held-out split tests temporal generalization within the dataset, not
a fresh sample. Out-of-foundation validation remains future work
(Chapter~\ref{chap:conclusion}).

\section{Conclusion Validity}
\label{sec:validity-conclusion}

Most metrics are point estimates from one time-aware held-out split. The
active-at-\(t\) margin additionally has a \(95\%\) repository-clustered bootstrap
interval excluding zero (Chapter~\ref{chap:evaluation}), and
Section~\ref{sec:results-label-sensitivity} tests label variants. Precision
remains conditional on sampled prevalence.

ROC-AUC is the primary ranking metric and is insensitive to class prevalence.
Average precision is additionally reported for the early-warning subset, whose
inactive rate is \(0.225\); it focuses on the positive class and uses that rate
as its no-skill reference \cite{saito2015precision}. It supports the same
conclusion while showing a larger margin over the trivial rules
(Section~\ref{sec:results-precision-recall}).
Two limitations remain. No precision-recall \emph{curve} is plotted, so the
shape of the trade-off at a given recall is not shown, only its area. And
precision and recall at a single operating point are still reported only at the
F1-selected threshold, which is chosen on validation data and is not optimal for
any particular triage budget. Qualitative case review is formative support, not
confirmatory evidence.
